\scanpage{15}
\mspage{11}
% source: top, continuation from scan 14
par ``effacement des sommets redondants'', \deleted{et}\unsure{ou} encore par ``composition des \deleted{\unsure{segments}} intervalles adjacents aux sommets redondants''. On a
\[
({\unsure{\Phi|\varphi}}^{\unsure{**|\circ *}})_{\unsure{0|o}}=\Phi_0\setminus\Phi_{\mathrm{red}}
\]
\[
\begin{aligned}
(\Phi_{\unsure{c|0|\epsilon|\cdot}}^{**})_1&=\deleted{\text{composé}}\text{ ens des intervalles}\\
&\qquad\deleted{\text{composés}}\inserted{\text{obtenus}}\text{ par composition des}\\
&\qquad\deleted{\text{\unsure{suite}}}\text{drapeaux maximaux}\\
&\qquad\uline{\text{connexes}}\text{ contenus dans }\Phi,
\end{aligned}
\]
% small c-like mark at lower right of Phi in the second display retained above; its status is uncertain
% source: middle, chain and relations
i.e. de la forme
\[
X_{\unsure{1|l}},x_1,X_2,x_2,\cdots,x_{\unsure{k|n}-1},X_k
\]
avec
\[
\deleted{x_1\triangleleft X_1,X_2,\ \unsure{x}}
\]
\[
\begin{aligned}
&X_1<x_1<X_2<x_2\ \cdots\ <x_{\unsure{n|k|u}-1}<X_n\\
&\text{et }x_1\triangleleft X_1,X_2,\quad x_2\unsure{\mid|,}\triangleleft X_2,X_3,\quad\cdots,\quad x_{n-1}\triangleleft X_{n-1},X_n.
\end{aligned}
\]

% source: middle-lower, components
\unsure{Elles|Ils} sont en corr. 1-1 avec les composantes connexes de $\Phi$, lesquelles n'en diffèrent que par l'adjonction éventuelle d'un $x_0\triangleleft X_1$, ou d'un $x_n\triangleleft X_n$, ou des deux.

% source: lower, shadow of a polygon
\deleted{\unsure{Et}} Ceci posé, à chaque \uline{polygone} $\Phi$, associons
\[
\operatorname{Omb}^{\unsure{\circ|0}}_{\unsure{M|\mathscr M}}(\Phi)=\deleted{\unsure{\forall X\in}}\bigcup_{X\in\Phi}\operatorname{Omb}^{\circ}_M(X)
=\left\{Z\in M\left|\begin{aligned}&\exists X\in\Phi,\\&Z\overset{\circ}{\ll}X\end{aligned}\right.\right\}
\]

% source: bottom, introduction to following proposition
Ceci posé, le résultat \unsure{\uline{clef}|clef} est le suivant
% archive sheet number: 14
